\section{Discussion}

\method{} does not prove that deleted data has no possible remaining effect. Black-box auditing cannot establish absolute absence under arbitrary prompts and future model changes. Instead, the framework provides bounded, reproducible evidence that $S_-$ is behaviorally close to $S_{\mathrm{clean}}$ under declared probe classes, adversary models, component interventions, and risk tolerances. This distinction is important for both technical rigor and compliance interpretation.

The main AI-methodology lesson is that deletion verification must follow data lineage while measuring counterfactual behavior. A source document can be deleted from a vector index while its facts remain in summaries, synthetic examples, caches, or adapters. Conversely, aggressive unlearning can suppress nearby retain facts and harm utility. A useful trustworthy-AI audit must therefore measure both deletion risk and retain utility, then use interventions to localize where the residual distinguishability is mediated.

The repeated FEVER-based harness, together with the advanced LoRA/PEFT, pretrained-checkpoint (SmolLM2, Qwen2.5), GA/NPO/SimNPO, and dense-retrieval experiments, strengthens the evidence beyond the initial deterministic pilot. Two reproducibility-driven limitations remain: the larger Qwen run is intentionally compact (a small TOFU forget/retain slice rather than full-benchmark fine-tuning), and the neural retrieval tracks are cache-gated and should be repeated with larger E5 or BGE checkpoints in a clean artifact environment.

\textbf{Threats to validity.}
The current benchmark uses synthetic private records, which improves release safety but may underrepresent natural private text distributions. The scorer validation table is computed on a released held-out labeled diagnostic set, complementing the human adjudication on redacted natural-data outputs. On that set the lexical localizers are deliberately conservative and imprecise (direct- and paraphrase-channel precision near $0.50$ and $0.25$), so the channels should be read as localizers rather than precise detectors; the formal pass/fail decision instead uses the two-sample counterfactual statistic, whose channel attains precision and recall of $1.0$ on the same set. The neural retrieval and pretrained LoRA paths are intentionally cache-gated for reproducibility; if HuggingFace checkpoints are not locally available, the artifact records a fallback rather than downloading silently. Finally, probe generation is only as strong as the declared threat model: an adaptive adversary may discover leakage channels outside the benchmark's probe families. Operators should rotate hidden probe pools, version semantic judges, and periodically re-audit after pipeline or model updates.

Several reviewer-relevant limitations remain explicit. First, configured synthetic failures are useful for regression testing; to reduce circularity, we now separately report deletion-operation labels that are independent of ADCU channel scores, the clean-counterfactual statistic, and human-adjudicated redacted cases. Second, the formal false-pass bound assumes independent probes from the declared audit distribution, whereas real probe families are correlated; the seed-block bootstrap and jitter tests are diagnostics, not substitutes for a full dependent-sampling theory. Third, the advanced PEFT/Qwen rows validate feasibility on cached public checkpoints but are too small to support high-precision confidence intervals, so tiny-$N$ intervals are treated as descriptive rather than strong statistical evidence. Fourth, the current membership, extraction, and influence-proxy baselines are audit-channel baselines, not complete reproductions of every published attack or benchmark evaluator.

\textbf{Pure black-box operation.}
When retrieval metadata is unavailable, \method{} still audits answer-side distinguishability and direct, paraphrase, extraction, and watermark evidence; the answer-only experiment removes artifact identifiers so retrieval dependence is credited only when it manifests in text. This is weaker than gray-box auditing but assumes no privileged access. If $S_{\mathrm{clean}}$ cannot be queried directly, the audit reports an approximation gap and uses the clean simulator, a filtered-retraining slice, or an archived clean index.

\textbf{Formal guarantees and AI assurance.}
\method{} complements rather than replaces certified removal, deletion-as-control, differential privacy, or pan-private continual release: it audits the deployed behavior, derivative artifacts, caches, summaries, adapters, and retrieval-control layers that sit outside the formal guarantee boundary, treating any component-level guarantee as prior evidence and spending more probes on uncovered downstream artifacts. It thus serves as an assurance layer that documents the clean reference, adversary class, probe coverage, allocation, interventions, confidence obtained, and where escalation remains necessary.
